\section{The Shape of the Cross}

Love of Christ must eventually meet suffering, but it must first refuse a lie: suffering is not good in itself. One suffers because some good is absent, damaged, threatened, or taken. Christianity proclaims the goodness of being; it does not call ruin holy merely because ruin hurts.\endnote{John Paul II, \emph{Salvifici Doloris} (11 February 1984), nos.~7--13: suffering is the experience of evil.}

Nor may suffering be read presumptively as punishment for a particular fault. God rebukes Job's friends: ``You have not spoken of me what is right.'' Of the man born blind, Christ says, ``Neither this man nor his parents sinned.''\endnote{Job 42:7; John 9:1--3. Divine judgment remains true; affliction alone does not disclose a particular guilt.} Known sin should be repented; affliction is not proof of what the sin was.

Whatever suffering can rightly be removed should be resisted by proportionate means. The Samaritan ``bound up his wounds'' and ``took care of him.'' Love does not pass by structural cruelty, preventable disease, abuse, hunger, loneliness, or unjust law and call inaction ``offering it up.'' John Paul II says the Gospel negates passivity: compassion must become effective help.\endnote{Luke 10:25--37; \emph{Salvifici Doloris}, nos.~28--30, including personal, professional, institutional, and social action.}

Some suffering nevertheless cannot now be removed. It may be endured without being minimized, advertised as a spiritual credential, or made the center of identity. Seeking medical care, reporting abuse, asking for help, lamenting, and naming evil truthfully do not betray the Cross. Christ himself says, ``My soul is very sorrowful, even to death,'' and asks that the cup pass. His final conformity is not numbness: ``Yet not what I will, but what you will.''

When unavoidable suffering is united to Christ, its evil is not renamed good. Love gives the sufferer's act a new relation to Christ's self-offering and Body. Paul says, ``I rejoice in my sufferings for your sake,'' and, ``In my flesh I complete what is lacking in Christ's afflictions for the sake of his body, that is, the Church.'' Nothing is lacking in the worth of Christ's Redemption; what opens is the redeemed person's participation in it for the members of Christ.\endnote{Col 1:24; \emph{Salvifici Doloris}, nos.~19--24: participation adds nothing to Christ's infinite redemptive worth.}

Because such participation is a free act of charity within grace, it can also be meritorious. Pain merits nothing by itself, and no creature earns the first grace or adds to the Cross. Yet, by Christ's gift, the sufferer's act may be offered for the graces needed for the sanctification and salvation of others---and, with charity's universal breadth, for all---without controlling another's freedom or guaranteeing a temporal outcome.\endnote{CCC 2007--2011, especially 2010, grounds every Christian merit in God's free initiative and teaches that, moved by the Holy Spirit and charity, one may merit for oneself and for others graces ordered to sanctification and eternal life.}

This possibility never licenses self-harm or imposed suffering. Voluntary penance belongs under prudence, moderation, state in life, health, duties, and competent spiritual direction. Severe asceticism belongs, if at all, only to cloistered contemplative or closely analogous vocations under competent direction; it cannot be generalized. Penance that damages bodily function or disables duty contradicts the good it claims to serve.\endnote{A boundary, not an ascetical prescription: fasting and bodily or devotional burdens are not proofs of greater love.}

Most participation in Christ will be less dramatic. Work honestly done, a child's need answered at an inconvenient hour, a grievance refused the luxury of hatred, a neighbor accompanied through illness, a sorrow borne without making others pay for it: these can become embodied prayer when charity gives them to God. Love of neighbor is not a diversion from union. When one loves the Beloved, one learns to love what the Beloved loves.

The Cross does not teach that pain is the end. It reveals what love can do when pain cannot yet be escaped. Resurrection remains the end.
